\subsection{Глубокое обучение с подкреплением на основе аппроксимации Q-функции}

\subsubsection{Постановка задачи и мотивация}

При реализации базовой модели алгоритма глубокого обучения с подкреплением мы опирались на работу \cite{mnih2015humanlevel} коллектива DeepMind. В ней решалась важная для современных когнитивных агентов задача -- переход от работы в простых детерминированных средах к обучению сложным стратегиям управления напрямую из высокоразмерных данных.

До появления этой архитектуры применимость RL-алгоритмов была сильно ограничена размерностью пространства состояний. При усложнении сред возникала проблема, заключающаяся в экспоненциальном росте объема данных от количества входных признаков, которую традиционно пытались обойти путем ручного определения признаков. Попытки напрямую скрестить глубокие нейросети и Q-обучение долгое время заканчивались неудачей. Также можно было наблюдать, как сильная зависимость по времени последовательных данных и постоянное изменение их распределения приводили к тому, что градиентный спуск делал процесс обучения крайне нестабильным, а алгоритм быстро расходился.

\subsubsection{Основные идеи и архитектура метода}

В алгоритме Deep Q-Network применяется нейронная сеть для аппроксимации функции оптимальной ценности действий $Q^∗(s,a)$. Стабилизировать обучение помогли два очень важных вспомогательных механизма. Первым таким решением стал отказ от поточного обучения в пользу механизма воспроизведения опыта. Вместо того чтобы обновлять веса сети на основе последних шагов, которые очень зависят друг от друга, агент непрерывно сохраняет свое текущее состояние, действие, награду и следующее состояние в специальный кольцевой буфер памяти. Обучение происходит на случайных мини-батчах, извлекаемых из этого буфера.

Вторым критическим компонентом стала архитектура с использованием изолированной целевой сети. На практике, если вычислять функцию потерь и одновременно обновлять веса одной и той же сети, возникает ситуация, когда целевые значения постоянно смещаются, что приводит к некорректным градиентам. Для решения этой проблемы добавлена фиксация целевой сети, из-за чего веса основной обучаемой модели обновляются на каждом шаге взаимодействия со средой, тогда как целевая сеть фиксируется и копируется из основной лишь раз в $C$ шагов.

\subsubsection{Результаты и выявленные ограничения}

Алгоритм DQN был протестирован на 49 играх платформы Atari 2600. Основным результатом стало то, что одна и та же архитектура нейросети с неизменными гиперпараметрами смогла достичь или превзойти уровень профессионального игрока-человека в большинстве игр.

Несмотря на это, глубокое обучение с подкреплением на базе DQN имеет ряд ограничений:

\begin{itemize}

  \item Процесс крайне чувствителен к шагу обучения и параметрам исследования среды.

  \item Агенту требуются миллионы тактов симуляции для схождения к оптимальной политике, так как обучение "с нуля" через градиентный спуск идет крайне медленно на ранних этапах.

\end{itemize}

Не стоит забывать и об ограничениях, оговоренных ранее.

\subsubsection{Связь с задачей моделирования когнитивных агентов}

Cтатья \cite{mnih2015humanlevel} демонстрирует, как классический подход решает задачу моделирования когнитивного агента через градиентную аппроксимацию Q-функции при фиксированной топологии мозга нейросети.

Анализ ограничений DQN формирует четкое обоснование исследовательских задач, поставленных в п. \ref{def:complexity}. В то время как классический RL тратит огромные вычислительные ресурсы на обновление весов в избыточной, вручную заданной нейросети, алгоритмы семейства NEAT предлагают альтернативный путь. NEAT решает проблему жесткой архитектуры, начиная с минимально возможной топологии и инкрементально усложняя ее только там, где это необходимо для выживания агента. Сравнение метрик (скорости схождения, количества параметров и устойчивости) между подходом, вдохновленным DQN, и алгоритмом NEAT является центральным элементом практической части данной работы.